% ============================================================================
\section{Discussion}
\label{sec:discussion}
% ============================================================================

The transverse rate a multi-echo spin echo measures is bulk relaxation plus an isotropic,
TE-linear surface rate $\rho\,(S/V)$ inseparable from it, with $S/V$ unobserved. The
consequence for the myelin water fraction is specific: surface relaxivity shortens the
intra/extra apparent $T_2$, and for the thinnest axons it drags that $T_2$ below the
myelin window, where the water is counted as myelin --- so fine white matter reads
myelin-\emph{richer} at fixed myelin content, by a substrate-dependent amount (${\sim}0.1$
percentage points at the cited $\rho$, super-linearly more --- $>1$\,pp --- at the upper
end of its order-of-magnitude uncertainty). It is no artefact: a wall-counting Monte Carlo
reproduces the rate from first principles, and the short-$T_2$ crossing is a property of
the window's own definition, present before any estimator sees the signal.

\paragraph{The same physics from two communities.} The contribution is partly one of
perspective: the time-dependent extra-axonal $D(t)$ of randomly packed media is well
established in diffusion MRI \citep{Novikov2014, Burcaw2015, Fieremans2016}, and the same
wall collisions, seen from relaxometry, are surface-relaxation events (Sec.~\ref{sec:theory:dual}).
The corollary: the ``$T_2$'' a relaxometry experiment reports is contaminated by a
geometry-dependent surface rate with no gradient played.

\paragraph{Implications for myelin water imaging.} MWF is read as a myelin-content index
and compared across tracts, subjects, and time. Because per-voxel MWF is noisy, studies
average it over a region of interest --- a callosal subregion, a tract, lobar white
matter --- and compare ROI means; subtracting two such means is where the surface term
bites. Each ROI carries its own offset set by its own thin-axon population, so
\begin{equation}
  \Delta\mathrm{MWF}_{\mathrm{app}}
  = \underbrace{\Delta\mathrm{MWF}_{\mathrm{true}}}_{\text{myelin difference}}
  \;+\;\big[\,\beta_A-\beta_B\,\big],
  \qquad \beta\ =\ F_I\,\Phi(\rho),
  \label{eq:roi_diff}
\end{equation}
where $F_I$ is the intra-axonal signal fraction and $\Phi(\rho)$ the volume-weighted
fraction of that intra-axonal water dragged below the myelin window (the thin-calibre tail),
so $\beta$ is larger for the finer region. The second bracket ---
a pure calibre difference --- adds to the myelin difference and cannot be separated from it
within the multi-echo data. It vanishes only when the two ROIs share their fine-axon
population, and is largest where calibre differs systematically: between callosal
subregions (the finer genu and splenium versus the large-diameter mid-body
\citep{Aboitiz1992} --- predicting the finer genu/splenium read myelin-\emph{richer} than
the mid-body at equal myelin), across development or ageing, or between patient and control
groups when disease shifts the calibre spectrum. At ${\sim}0.1$\,pp (${\sim}0.24$\,pp of
inferred myelin content) at the cited $\rho$ it is a small fraction of a gross regional
contrast, but a meaningful one for the subtle, ${\sim}1$\,pp differences cross-sectional
myelin studies pursue --- and, being deterministic in the substrate, a cohort-level
systematic that shifts a group mean rather than averaging out. The comparison it spares is
the field's cleanest, the \emph{same region of the same subject over time}: there the offset
is a nearly-constant, lumen-set term that cancels in the change
(Sec.~\ref{sec:results:demyel}), because primary demyelination leaves the inner diameters
the crossing depends on intact. None of this invalidates MWF; it bounds a contribution of
the same $S/V$-driven family as the documented exchange bias \citep{Harkins2012}, not
previously quantified as a surface-relaxivity term inseparable from the bulk $T_2$.

\paragraph{What counts as myelin water.} The mechanism is not really about axons: the
short-$T_2$ window captures \emph{any} water whose $T_2$ is dragged below it by surface
relaxivity, which happens once $S/V$ exceeds
$(1/T_2^{\mathrm{cut}}-1/T_{2,\mathrm b})/\rho\approx 23\,\mu$m$^{-1}$ at the cited $\rho$
(a sub-micron scale) and $\approx 11\,\mu$m$^{-1}$ at the upper $\rho$. This is the very
physics that makes bilayer-trapped myelin water short-$T_2$ in the first place, so MWF has
always measured ``high-$S/V$ trapped water'', of which myelin is merely the dominant source
in health. It follows --- mechanism-predicted, not simulated here --- that \emph{any}
membrane-dense sub-micron compartment qualifies: myelin debris, fine reactive-glial
processes or microvacuolation could add short-$T_2$ signal counted as myelin and prop up
apparent MWF, while whole cells ($S/V\!<\!1\,\mu$m$^{-1}$) are far too large to cross. This
widens the caveat from a calibre confound to a statement about what the short-$T_2$ pool
contains, and is a natural target for a drop-in Monte Carlo of explicit pathological
microstructure.

\paragraph{Idealised packing, and where Monte Carlo becomes essential.} Our closed forms
are exact for the substrate we assume --- parallel, circular, non-touching cylinders packed
by random sequential adsorption --- and for that geometry they are also the \emph{efficient}
route. A Monte-Carlo pack pays a sub-step cost that scales as $1/R^2$ in its smallest axon
(an impermeable-wall relaxivity weight needs steps $\ll R$), so the very calibres that
dominate the bias are the most expensive to simulate \citep{Hall2009}, while the analytics
returns them for free. The idealisation has a specific blind spot,
though. The exterior rate uses the \emph{mean} $S/V$ --- total wall perimeter over total
extra-axonal area --- and a mean cannot see the \emph{distribution of local} $S/V$ that
realistic packing creates: where axons deform, touch, undulate or cross \citep{Xu2018},
extra-axonal water is trapped in thin interstitial crevices of very high local $S/V$, and
the window-crossing the intra-axonal calibre tail undergoes could occur there too ---
extra-axonal water read as myelin. Whether it does hinges on motional averaging, and two
\emph{distinct} length scales must not be conflated. A local sheet of high $S/V$ shortens the
$T_2$ enough to cross the myelin window once its width falls below
${\sim}90$\,nm ($S/V\gtrsim23\,\mu$m$^{-1}$ at the cited $\rho$) --- a scale the extra-axonal
space genuinely reaches. But whether that water \emph{averages} with the bulk is a
\emph{separate} condition, set by how fast a spin escapes the pocket: at a $90$\,nm width the
traversal time $w^2/D\!\sim\!4\times10^{-3}$\,ms is far below the echo spacing, so an
\emph{open} $90$\,nm sheet averages comfortably and feels only the mean $S/V$. Averaging fails
only where the pocket is kinetically \emph{isolated} --- narrow throats or near-contacts whose
bottleneck traversal time approaches the echo spacing --- a percolation-scale property of the
packing, not of the sheet width, and one we do not quantify here. So the mean-field rate is a
lower bound (Jensen; Sec.~\ref{sec:theory:closedforms}) and the size of the correction is set by
this unquantified isolation scale. This is the regime the analytics cannot reach and only a
Monte Carlo of non-idealised geometry can resolve --- and, by the same $1/R^2$ argument, the
most costly to run. The reading is therefore a duality, not a ranking: the closed form should carry the
idealised packing, where Monte Carlo is wasteful, and Monte Carlo should be reserved for the
irreducibly non-idealised part --- crevices, undulation, crossings --- where the closed form
has nothing to say.

\paragraph{Why the correction injects the fraction rather than fitting it.} The quantity a
diffusion model returns is the intra-axonal \emph{signal} fraction, not the geometric fibre
volume fraction $VF$; a spherical-mean fit even ties its extra-axonal (tortuosity) diffusivity to
that signal fraction, so the packing geometry that sets $S/V$ is never an independent output.
Because the surface term is $b$-independent and inseparable from the bulk $T_2$, $S/V$ --- and
hence $VF$ --- is not identifiable from the fit, and feeding the fitted fraction back into the
surface term would be circular (it is the very quantity the surface term biases). The correction
therefore treats $VF$ as a \emph{bounded prior} supplied from regional histology --- exactly the
ingredient a cross-TE cohort harmonisation needs and can plausibly obtain --- and leaves the fit
to return only the corrected fraction.

\paragraph{What could resolve it, and what cannot.} The degeneracy is structural, so no
refinement of the multi-echo acquisition alone breaks it: echo spacing, train length and
signal-to-noise change precision, not the non-identifiability. Diffusion is, in
principle, the orthogonal axis --- the same wall collisions, read with a gradient. The
elegant point is that the disorder enters diffusion through the \emph{same} fibre
orientation distribution $\mathcal F(\hat{\mathbf n})$ as the signal convolution
Eq.~\eqref{eq:fod_conv}: in the low-$b$, short-time regime the apparent diffusivity along
$\hat{\mathbf g}$ is that FOD convolved with a disorder kernel,
\begin{equation}
  \frac{D(t;\hat{\mathbf g})}{D_0}=\int_{S^2}\mathcal F(\hat{\mathbf n})\,
  \kappa_{\mathrm{dis}}(\hat{\mathbf g}\!\cdot\!\hat{\mathbf n};t)\,\mathrm d\hat{\mathbf n},
  \quad
  \kappa_{\mathrm{dis}}=1-\tfrac{4}{3\sqrt\pi}\sqrt{D_0 t}\,
  \Big[\textstyle\sum_c f_c\,(S/V)_c\Big]\big[1-(\hat{\mathbf g}\!\cdot\!\hat{\mathbf n})^2\big],
  \label{eq:kdis}
\end{equation}
with $(S/V)_c$ the interior $\svin$ or exterior $S_{\mathrm{ext}}/V_{\mathrm{ext}}$. Because
$1-(\hat{\mathbf g}\!\cdot\!\hat{\mathbf n})^2=\tfrac23[1-P_2(\hat{\mathbf g}\!\cdot\!\hat{\mathbf n})]$
carries only $\ell=0$ and $\ell=2$, the surface disorder is a \emph{rank-2} function on the
sphere --- round ($\ell{=}0$) to relaxation, fibre-shaped ($\ell{=}0{+}2$) to diffusion --- and
spherical convolution preserves that band limit for \emph{any} FOD. The orientation-invariant
total $S/V$ is then the $\ell=0$ (trace) part, which an isotropic (trace) oscillating-gradient
acquisition or a spectrally tuned $b$-tensor \citep{Lundell2019, Westin2016} recovers
independent of fibre dispersion. In practice it is gradient-limited: the short-time regime needs
OGSE frequencies whose deliverable $b$-value collapses as $\sim$$1/f^3$ \citep{Does2003} ---
sitting at the corner frequency of a micron axon needs $G\!\sim\!70$--$260\,$T/m, one to two
orders beyond even preclinical inserts --- so the fine calibres that carry the effect lie out of
reach, and a realisable protocol is left to future work. What no route removes
is the absolute calibration: the inner- and outer-wall relaxivities
$\rho_{\mathrm i},\rho_{\mathrm e}$ --- which may differ, and scale the bias --- are not
independently established for myelinated tissue, and the myelin-water pool itself, with its
very large bilayer $S/V$, is plausibly surface-dominated in a way we have not modelled
here.

\paragraph{Relation to exchange and other MWF confounds.} MWF is already known to
depend on factors other than myelin content: $B_1$/flip-angle imperfections, which
the extended-phase-graph fit corrects \citep{Prasloski2012, Lebel2010}; iron and
susceptibility; and inter-compartment water \emph{exchange}, which Harkins, Dula and
Does showed makes the apparent MWF depend on axon size \citep{Harkins2012, Dula2010}.
The exchange bias is, like ours, driven by surface-to-volume ratio, and both are largest
at high $S/V$. They are \emph{mechanistically distinct} --- exchange moves magnetisation
between pools of different $T_2$, whereas surface relaxivity destroys coherence at the wall
and, for the thinnest axons, moves intra water across the window --- but \emph{operationally
entangled}: both act at the small-calibre end and neither is separable within the multi-echo
data. We deliberately isolate surface relaxivity by making the myelin impermeable in the
simulation, so the rate validated in Fig.~\ref{fig:signal}A and the bias of
Fig.~\ref{fig:mwf} are exchange-free; we do not claim to have measured their \emph{relative}
magnitude, which would require permeable walls and is left to future work. Myelin is a strong
barrier (slow apparent exchange, $1$--$10\,$s$^{-1}$, permeation at nodes and paranodes
\citep{Nilsson2013}), whereas wall \emph{encounters} recur on the perpendicular diffusion
correlation time ($\sim$$0.1$--$1\,$ms): a spin bounces two to three orders of magnitude more
often than it permeates, which is what makes isolating surface relaxivity with $\kappa=0$
physically justified.

\paragraph{Biological idealisations.} Four substrate assumptions bound the biological
reading. First, the effect is carried by the sub-micron calibre tail (the crossing is at
$d^{\ast}\!\approx\!0.17\,\mu$m at the cited $\rho$), so the magnitude depends on how much
sub-micron water the tissue holds --- which we anchor with rat spinal-cord cross-sections
\citep{Zaimi2018}, a coarser-calibre tissue than the human cortical and callosal white matter
that is the clinical target. Human callosal fibres run finer, with a substantial sub-micron
population \citep{Aboitiz1992}, so if anything the rat-cord anchor \emph{under}-states the
tissue-borne effect; a human-EM calibre census is the proper anchor and we do not claim one.
Second, we hold the g-ratio fixed at $0.70$ across calibre, whereas fine axons are relatively
more thickly myelinated (lower $g$) \citep{Stikov2015}; since the crossover
$\mathrm{VF}^{\ast}=1/(1+g)$ depends on $g$, the single sign law is a calibre-averaged
statement, and a realistic $g(d)$ would spread the crossover across calibres rather than
abolish it. Third, we assign the inner and outer walls one relaxivity $\rho_{\mathrm i}=\rho_{\mathrm e}$;
these are chemically different membranes whose relaxivities may differ, a \emph{structural}
(not merely scale) uncertainty that could shift the interior/exterior balance and hence the
$f_{\mathrm{intra}}$ sign, and $\rho$ itself is imported from a histology-fit model coefficient
\citep{Barakovic2023} that may partly absorb exchange or susceptibility. Fourth, the
longitudinal cancellation assumes lumen-preserving (pure g-ratio) demyelination; real MS and
Wallerian degeneration co-vary axon calibre through swelling, beading, atrophy and outright
axonal loss, which alter the very inner-diameter distribution the offset is set by --- so the
clean within-subject cancellation is an idealised bound, and where pathology shifts the calibre
spectrum the cross-region confound re-enters the longitudinal comparison.

\paragraph{Limitations.} We assumed ideal, instantaneous refocusing, isolating the
bulk-versus-surface question; finite-pulse and susceptibility effects are real but
orthogonal and treated elsewhere. This sets the effect in perspective: the
flip-angle/stimulated-echo and fibre-orientation/susceptibility \citep{Birkl2021}
systematics the MWF field already corrects for are themselves of order a percentage point
or more --- comparable to or larger than the surface bias at the cited $\rho$, which is a
further, hitherto-unnamed term of the same family. The MWF magnitudes come from a forward
model carrying the validated rate over a calibre-distributed intra pool
(Sec.~\ref{sec:methods:mwf}), so the absolute numbers inherit its two-/three-pool,
non-exchanging, canonical-Gamma idealisations --- which modulate the magnitude but not the
existence or sign of the term. The robust core needs no estimator: the closed-form rate and
its wall-counting Monte-Carlo validation (a motional-narrowing limit and a geometry-only
simulation with no relaxivity model) agree, and the short-$T_2$ crossing they imply is a
grid-free property of the signal. The claim we stand behind is the inseparability itself,
the sign of its consequence (fine reads myelin-richer), and the order of magnitude that
follows once $\rho$ is fixed.
